% !TEX root = 00_instructivo.tex
\chapter{Transductores mecánicos: Acelerómetro y podómetro}

\section*{Objetivos}
\begin{itemize}
    \item Comprender el principio de funcionamiento de un acelerómetro MEMS (MPU6050).
    \item Implementar un algoritmo de conteo de pasos (podómetro) usando la magnitud del vector de aceleración.
    \item Evaluar la precisión del podómetro en diferentes condiciones de movimiento (caminar, correr, movimientos bruscos).
    \item Alimentar el sistema con una batería portable para demostrar su uso en aplicaciones de monitoreo de actividad.
\end{itemize}

\section*{Preguntas previas}
\begin{enumerate}
    \item ¿Qué es una IMU (Unidad de Medición Inercial) y qué grados de libertad ofrece el MPU6050?
    \item ¿Por qué se usa la magnitud de la aceleración (\( \sqrt{a_x^2 + a_y^2 + a_z^2} \)) y no un solo eje para detectar pasos?
    \item ¿Qué es el \textit{efecto de histéresis} en el contexto de un umbral de detección?
    \item Investigue: ¿qué es el DMP (Digital Motion Processor) del MPU6050 y para qué sirve?
\end{enumerate}

\section*{Materiales y equipo}
\begin{itemize}
    \item Arduino Uno (o Arduino Nano para versión portable).
    \item Módulo MPU6050 (acelerómetro + giróscopo 6-DOF).
    \item Batería recargable de 9V o power bank 5V (para alimentación portable).
    \item Conector para batería (clips o cable USB).
    \item Mini protoboard.
    \item Cables Dupont.
    \item Cinta o velcro para fijar el conjunto a la cintura o al tobillo.
    \item Computadora con IDE Arduino (para carga inicial y monitoreo).
\end{itemize}

\section*{Procedimiento}
\begin{enumerate}
    \item \textbf{Conexión y prueba del MPU6050:}
    \begin{enumerate}
        \item Conecte el MPU6050 al Arduino (VCC a 5V, GND a GND, SDA a A4, SCL a A5).
        \item Cargue el código de la Sección 4.1 de la guía (podómetro).
        \item Verifique en el monitor serial que se leen valores de aceleración.
    \end{enumerate}
    \item \textbf{Calibración del umbral:}
    \begin{enumerate}
        \item Sujete el MPU6050 firmemente a su cintura o tobillo.
        \item Camine 20 pasos en un lugar plano mientras observa el monitor serial.
        \item Ajuste la constante \texttt{THRESH} en el código (por ejemplo, 1.2, 1.5, 1.8) hasta que el conteo sea preciso (error < 2 pasos).
    \end{enumerate}
    \item \textbf{Prueba en diferentes escenarios:}
    \begin{enumerate}
        \item Realice 3 pruebas de 50 pasos cada una: caminando normalmente, caminando rápido y subiendo escaleras.
        \item Anote los pasos contados por el Arduino y compárelos con los pasos reales.
        \item Calcule el error porcentual para cada escenario.
    \end{enumerate}
    \item \textbf{Modo portable:}
    \begin{enumerate}
        \item Desconecte el Arduino de la computadora y conéctelo a la batería portable.
        \item (Opcional) Agregue una pantalla LCD de 16x2 para mostrar el conteo en tiempo real.
        \item Camine 100 pasos en un circuito cerrado y al final reconecte a la PC para leer el conteo.
    \end{enumerate}
\end{enumerate}

\section*{Esquemático}
\begin{figure}[htbp]
\centering
\begin{circuitikz}[scale=0.8, transform shape]
\draw (0,0) node[above]{Arduino Nano (portable)};
\draw (1,0) node[ground]{GND};
\draw (1,0.5) node[anchor=south]{5V};
\draw (0.5,1) node[anchor=east]{A4 (SDA)} -- (2.5,1);
\draw (0.5,1.5) node[anchor=east]{A5 (SCL)} -- (2.5,1.5);
\node at (3,1.25) {MPU6050};
% Batería
\draw (1,2.5) node[anchor=west]{Batería 9V} -- (2.5,2.5);
\draw (1,2.5) to[short, *-] (1,2);
\draw (1,2) to[short] (1,0.5);
\draw (2.5,2.5) node[ground]{GND};
\end{circuitikz}
\caption{Conexión del MPU6050 a un Arduino Nano alimentado por batería. El bus I2C es compartido.}
\end{figure}

\section*{Preguntas de análisis}
\begin{enumerate}
    \item ¿Por qué la magnitud de la aceleración es más robusta para la detección de pasos que usar solo el eje Z?
    \item ¿Qué limitaciones tiene este algoritmo podómetro? ¿Cómo se podría mejorar usando el giróscopo?
    \item ¿Por qué se incluye la variable \texttt{stepActive} y el umbral \texttt{THRESH * 0.5}? Explique su función para evitar doble conteo.
    \item Compare el consumo energético del sistema con y sin batería. ¿Es viable para un uso de 24 horas?
    \item Proponga una modificación al código para estimar la distancia recorrida (requiere conocer la longitud del paso).
\end{enumerate}

\section*{Bibliografía}
\begin{itemize}
    \item Hoja de datos MPU6050. InvenSense
    \item “Step counting using accelerometers” – Application Note, Analog Devices.
\end{itemize}
